% ---------------------------------------------------------------------------
% Shared preamble for the AIG230 final report.
% ---------------------------------------------------------------------------

\usepackage[utf8]{inputenc}
\usepackage[T1]{fontenc}
\usepackage[english]{babel}

% --- Typography ------------------------------------------------------------
% Libertinus is a Linux Libertine derivative: an old-style serif with a large
% x-height that stays readable at 11pt, paired with a matching sans for
% headings. Chosen to echo the typographic feel of distill.pub articles, whose
% figure style this report also follows.
\usepackage{libertinus}
\usepackage{microtype}

\usepackage[a4paper,
            top=2.5cm, bottom=2.5cm, left=2.7cm, right=2.7cm,
            headheight=14pt]{geometry}

% --- Mathematics -----------------------------------------------------------
\usepackage{amsmath}
\usepackage{amssymb}
\usepackage{bm}

% --- Tables and figures ----------------------------------------------------
\usepackage{booktabs}
\usepackage{longtable}
\usepackage{array}
\usepackage{multirow}
\usepackage{graphicx}
\usepackage{float}
\usepackage[font=small,labelfont=bf,justification=raggedright,
            singlelinecheck=false]{caption}
\usepackage{subcaption}

% --- Colour ----------------------------------------------------------------
\usepackage{xcolor}
\definecolor{accentblue}{HTML}{2A78D6}
\definecolor{accentorange}{HTML}{EB6834}
\definecolor{accentgreen}{HTML}{1BAF7A}
\definecolor{inkprimary}{HTML}{0B0B0B}
\definecolor{inksecondary}{HTML}{52514E}
\definecolor{inkmuted}{HTML}{898781}
\definecolor{rule}{HTML}{E1E0D9}
\definecolor{surface}{HTML}{FCFCFB}

% --- Boxes -----------------------------------------------------------------
\usepackage[most]{tcolorbox}

\newtcolorbox{keypoint}[1][]{
  enhanced, breakable,
  colback=surface, colframe=accentblue,
  boxrule=0pt, leftrule=2.5pt,
  arc=1pt, left=8pt, right=8pt, top=6pt, bottom=6pt,
  fontupper=\small, #1
}

\newtcolorbox{decisionbox}[1][]{
  enhanced, breakable,
  colback=surface, colframe=accentgreen,
  boxrule=0pt, leftrule=2.5pt,
  arc=1pt, left=8pt, right=8pt, top=6pt, bottom=6pt,
  fontupper=\small, #1
}

% Marks a paragraph whose wording depends on final trained numbers.
\newtcolorbox{pending}[1][]{
  enhanced, breakable,
  colback=yellow!6, colframe=accentorange,
  boxrule=0pt, leftrule=2.5pt,
  arc=1pt, left=8pt, right=8pt, top=6pt, bottom=6pt,
  fontupper=\small\itshape, #1
}

% --- Headings --------------------------------------------------------------
\usepackage{titlesec}
\titleformat{\section}{\Large\bfseries\color{inkprimary}}{\thesection}{0.7em}{}
\titleformat{\subsection}{\large\bfseries\color{inkprimary}}{\thesubsection}{0.6em}{}
\titleformat{\subsubsection}{\normalsize\bfseries\color{inksecondary}}
            {\thesubsubsection}{0.6em}{}
\titlespacing*{\section}{0pt}{2.0ex plus 1ex minus .2ex}{1.1ex plus .2ex}

% --- Lists -----------------------------------------------------------------
\usepackage{enumitem}
\setlist{itemsep=2pt, topsep=4pt, parsep=0pt}

% --- Headers ---------------------------------------------------------------
\usepackage{fancyhdr}
\pagestyle{fancy}
\fancyhf{}
\fancyhead[L]{\small\color{inkmuted}AIG230 Final Project --- Group 7}
\fancyhead[R]{\small\color{inkmuted}\thepage}
\renewcommand{\headrulewidth}{0.4pt}
\renewcommand{\headrule}{\hbox to\headwidth{\color{rule}\leaders\hrule height \headrulewidth\hfill}}

% --- Code ------------------------------------------------------------------
\usepackage{listings}
\lstset{
  basicstyle=\ttfamily\footnotesize,
  keywordstyle=\color{accentblue}\bfseries,
  commentstyle=\color{inkmuted}\itshape,
  stringstyle=\color{accentgreen},
  showstringspaces=false,
  breaklines=true,
  frame=leftline,
  framesep=6pt,
  rulecolor=\color{rule},
  xleftmargin=10pt,
  numbers=none,
  language=Python,
}

% --- Links -----------------------------------------------------------------
\usepackage[hidelinks,
            colorlinks=true,
            linkcolor=accentblue,
            citecolor=accentblue,
            urlcolor=accentblue]{hyperref}
\usepackage{url}

% --- Spacing ---------------------------------------------------------------
\setlength{\parskip}{0.55em}
\setlength{\parindent}{0pt}
\renewcommand{\arraystretch}{1.15}

% --- Convenience -----------------------------------------------------------
\newcommand{\code}[1]{\texttt{\small #1}}
\newcommand{\dmodel}{d_{\mathrm{model}}}
\newcommand{\dff}{d_{\mathrm{ff}}}
\newcommand{\dk}{d_k}
\newcommand{\entos}{EN\,$\rightarrow$\,ES}
\newcommand{\estoen}{ES\,$\rightarrow$\,EN}

% ---------------------------------------------------------------------------
% Fallbacks so the document compiles before any model has been trained.
% `generated/macros.tex` overrides whichever of these it defines.
% ---------------------------------------------------------------------------
\newcommand{\definefallbacks}{%
  \providecommand{\HasResults}{false}%
  \providecommand{\PrimaryRun}{BPE 16k, learned embeddings}%
  % corpus
  \providecommand{\TatoebaEnglishSentences}{--}%
  \providecommand{\TatoebaSpanishSentences}{--}%
  \providecommand{\TatoebaLinksScanned}{--}%
  \providecommand{\PairsRaw}{--}%
  \providecommand{\PairsClean}{--}%
  \providecommand{\RetentionPercent}{--}%
  \providecommand{\PairsTrain}{--}%
  \providecommand{\PairsValidation}{--}%
  \providecommand{\PairsTest}{--}%
  \providecommand{\Components}{--}%
  \providecommand{\Leakage}{--}%
  \providecommand{\LargestComponent}{--}%
  \providecommand{\TypesEn}{--}%
  \providecommand{\TypesEs}{--}%
  \providecommand{\TypeRatio}{--}%
  \providecommand{\TokensEn}{--}%
  \providecommand{\TokensEs}{--}%
  \providecommand{\HapaxEnPercent}{--}%
  \providecommand{\HapaxEsPercent}{--}%
  \providecommand{\OovEnPercent}{--}%
  \providecommand{\OovEsPercent}{--}%
  \providecommand{\FertilityEn}{--}%
  \providecommand{\FertilityEs}{--}%
  \providecommand{\SubwordVocab}{--}%
  \providecommand{\WordVocab}{--}%
  \providecommand{\MuseCoveragePercent}{--}%
  \providecommand{\MuseCovered}{--}%
  \providecommand{\MuseRandom}{--}%
  % per-run parameters, epochs and scores
  \providecommand{\ParamsBpeScratch}{--}%
  \providecommand{\ParamsWordRandom}{--}%
  \providecommand{\ParamsWordMuse}{--}%
  \providecommand{\ParamsLstmBaseline}{--}%
  \providecommand{\EpochsBpeScratch}{--}%
  \providecommand{\EpochsWordRandom}{--}%
  \providecommand{\EpochsWordMuse}{--}%
  \providecommand{\EpochsLstmBaseline}{--}%
  \providecommand{\BleuBpeScratchEnEs}{--}%
  \providecommand{\BleuBpeScratchEsEn}{--}%
  \providecommand{\BleuBpeScratchMean}{--}%
  \providecommand{\BleuWordRandomEnEs}{--}%
  \providecommand{\BleuWordRandomEsEn}{--}%
  \providecommand{\BleuWordRandomMean}{--}%
  \providecommand{\BleuWordMuseEnEs}{--}%
  \providecommand{\BleuWordMuseEsEn}{--}%
  \providecommand{\BleuWordMuseMean}{--}%
  \providecommand{\BleuLstmBaselineEnEs}{--}%
  \providecommand{\BleuLstmBaselineEsEn}{--}%
  \providecommand{\BleuLstmBaselineMean}{--}%
  \providecommand{\ChrfBpeScratchEnEs}{--}%
  \providecommand{\ChrfBpeScratchEsEn}{--}%
}

% Include a figure only if the file exists, so a partially-populated
% figures directory never breaks the build.
%
% The figure names produced by nmt.viz contain underscores (`data_construction`),
% which TeX reads as subscript markers and which would otherwise raise
% "Missing $ inserted" the moment the name reaches \includegraphics.
% \detokenize turns the argument into plain characters before that happens.
% \IfFileExists searches the current directory and TEXINPUTS -- it does *not*
% consult \graphicspath. Passing a bare name therefore always took the "not
% found" branch and silently rendered a placeholder box for every figure, even
% though the PDFs were present. The directory is given explicitly instead;
% each document sets \figuredir to point at reports/figures/.
\providecommand{\figuredir}{../figures/}

\newcommand{\maybefigure}[4]{%
  \IfFileExists{\figuredir\detokenize{#1}.pdf}{%
    \begin{figure}[htbp]
      \centering
      \includegraphics[width=#2\linewidth]{\figuredir\detokenize{#1}}
      \caption{#3}
      \label{#4}
    \end{figure}%
  }{%
    \begin{figure}[htbp]
      \centering
      \fbox{\parbox{0.8\linewidth}{\centering\small\itshape
        Figure not generated yet: \code{\detokenize{#1}.pdf}\\[2pt]
        Run \code{python -m nmt.viz.make\_figures}.}}
      \caption{#3}
      \label{#4}
    \end{figure}%
  }%
}
